\chapter{Container Workloads}
\label{ch:workloads}

This chapter demonstrates practical container workloads running on
Veilbox. These examples validate that the container runtime works
correctly and provide reference configurations for common use cases.

\section{Prerequisites}

All examples assume:
\begin{itemize}[nosep]
  \item Veilbox is booted and the root shell is available (via console
        or SSH).
  \item containerd is running (verified with \texttt{ps | grep containerd}).
  \item Networking is functional (verified with \texttt{ping 10.0.2.2}).
  \item The state partition is mounted at /mnt/state.
\end{itemize}

\section{Example 1: Hello World}

The simplest container test:

\begin{lstlisting}
nerdctl pull alpine:latest
nerdctl run alpine:latest echo "Hello from Veilbox!"
\end{lstlisting}

Expected output:
\begin{lstlisting}
Hello from Veilbox!
\end{lstlisting}

\section{Example 2: Interactive Shell}

Run an interactive Alpine shell:

\begin{lstlisting}
nerdctl run -it alpine:latest sh
/ # apk add --no-cache curl
/ # curl https://example.com
/ # exit
\end{lstlisting}

\section{Example 3: Web Server}

Run nginx with host port mapping:

\begin{lstlisting}
nerdctl run -d \
    --name web \
    -p 80:80 \
    nginx:alpine

# Test from inside Veilbox:
wget -q -O- http://localhost

# Test from host (via QEMU port forwarding):
curl http://localhost:8080/
\end{lstlisting}

\section{Example 4: Persistent Data}

Use a bind mount for persistent container data:

\begin{lstlisting}
# Create data directory on state partition
mkdir -p /mnt/state/data

# Run container with bind mount
nerdctl run -d \
    --name web-data \
    -v /mnt/state/data:/usr/share/nginx/html:ro \
    nginx:alpine

# Write data from host
echo "Persistent" > /mnt/state/data/index.html
\end{lstlisting}

\section{Example 5: Resource Limits}

Test resource constraints:

\begin{lstlisting}
# Memory limit
nerdctl run --memory=64m alpine sh -c \
    "dd if=/dev/zero of=/dev/null bs=1M count=100"

# CPU limit (0.5 cores)
nerdctl run --cpus=0.5 alpine sh -c \
    "dd if=/dev/zero of=/dev/null bs=1M count=1000"
\end{lstlisting}

\section{Example 6: Multi-Container}

Start multiple containers and verify isolation:

\begin{lstlisting}
# Start two HTTP servers
nerdctl run -d --name web1 nginx:alpine
nerdctl run -d --name web2 nginx:alpine

# They have separate network stacks
nerdctl exec web1 ip addr show eth0
nerdctl exec web2 ip addr show eth0

# They can communicate via the host bridge
nerdctl exec web1 ping -c 1 $(nerdctl inspect web2 | grep IPAddress)
\end{lstlisting}

\section{Troubleshooting Containers}

Common issues:
\begin{lstlisting}
nerdctl logs <container>     # View stdout/stderr
nerdctl inspect <container>  # Detailed status
nerdctl stats                # Resource usage
nerdctl top <container>      # Running processes
nerdctl events               # Stream container events
\end{lstlisting}


\section{Example 7: Database Container}

Run MariaDB with persistent storage:

\begin{lstlisting}
# Create data directory
mkdir -p /mnt/state/mariadb

# Run MariaDB container
nerdctl run -d \
    --name mariadb \
    -v /mnt/state/mariadb:/var/lib/mysql \
    -e MARIADB_ROOT_PASSWORD=example \
    -p 3306:3306 \
    mariadb:10.11

# Connect from another container
nerdctl run -it --rm mariadb:10.11 \
    mysql -h 10.0.2.15 -u root -pexample
\end{lstlisting}

\section{Example 8: Python Web App}

Deploy a Python Flask application:

\begin{lstlisting}
# Create app directory
mkdir -p /mnt/state/app

# Write a simple Flask app
cat > /mnt/state/app/app.py << 'EOF'
from flask import Flask
app = Flask(__name__)

@app.route('/')
def hello():
    return "Hello from Veilbox Container!"

if __name__ == '__main__':
    app.run(host='0.0.0.0', port=5000)
EOF

# Run with Python image
nerdctl run -d \
    --name flask-app \
    -v /mnt/state/app:/app \
    -p 5000:5000 \
    python:3.12-slim \
    python /app/app.py

# Test
wget -q -O- http://localhost:5000
\end{lstlisting}

\section{Example 9: Resource Monitoring}

Run a monitoring container with host metrics:

\begin{lstlisting}
# Run a stress test container
nerdctl run -d --name stress --memory=128m alpine \
    sh -c "while true; do dd if=/dev/zero of=/dev/null bs=1M count=100; done"

# Monitor from host
while true; do
    clear
    echo "=== Container Stats ==="
    nerdctl stats --no-stream
    echo "=== Host Memory ==="
    free -m
    echo "=== Host CPU ==="
    top -bn1 | head -5
    sleep 2
done
\end{lstlisting}

\section{Example 10: Network Test}

Test container networking with iperf:

\begin{lstlisting}
# Start iperf server
nerdctl run -d --name iperf-server -p 5201:5201 \
    alpine sh -c "apk add iperf3 && iperf3 -s"

# Run iperf client from another container
nerdctl run --rm alpine sh -c \
    "apk add iperf3 && iperf3 -c 10.0.2.15"
\end{lstlisting}

\section{Example 11: Cleanup}

Remove all containers and images:

\begin{lstlisting}
# Stop all running containers
nerdctl stop $(nerdctl ps -q) 2>/dev/null

# Remove all containers
nerdctl rm $(nerdctl ps -aq) 2>/dev/null

# Remove all images
nerdctl rmi $(nerdctl images -q) 2>/dev/null

# Clean containerd content store
rm -rf /mnt/state/containerd/*
\end{lstlisting}
